\uuid{L76N}
\titre{Matrice d'une application linéaire/Noyau/Image/Changement de bases}
\auteur{hanani ; mzari}
\organisation{exo7}
\licence{CC BY-NC-SA 4.0}
\datecreate{2020-01-20}
\niveau{L1}
\module{Algèbre}
\chapitre{Matrice}
\sousChapitre{Matrice et application linéaire}
\difficulte{3}
\theme{Matrice d'une application linéaire, Noyau, Image, Changement de bases}
\exo7qcmid{lille-2.matrice-application.020}

\contenu{
\begin{question}
\qclassification{108.03}
\qtags{motcle=Matrice d'une application linéaire/Noyau/Image/Changement de bases}
Soit $f$ l'endomorphisme de $\Rr^3$ dont la matrice dans la base canonique ${\cal B}=\{e_1,e_2,e_3\}$ est :
$$A = \left(\begin{array}{rcc}
-1&1&1\\0&1&1\\ 0&-2&4\end{array}\right).$$
On note $Id$ l'application identité de $\Rr^3$.
Quelles sont les assertions vraies ?
\begin{answers}
\good{$\dim\ker (f+Id)=\dim \ker (f-2Id)= \dim \ker (f-3Id)=1 $}.
\bad{$\dim\ker (f+Id)=\dim\ker (f-2Id)=1$ et $\dim\ker (f-3Id)=2$}.
\good{Il existe une base de $\Rr^3$ dans laquelle la matrice de $f$ est :
$B= \left(\begin{array}{rcc}-1&0&0\\
0&2&0\\ 0&0&3\\\end{array}\right)$}.
\good{L'application $(f+Id)o(f-2Id)o(f-3Id)$ est nulle}.
\end{answers}
\begin{explanations} On vérifie que $\{u_1=(1,0,0)\}$ est une base de $\ker (f+Id)$,
$\{u_2=(2,3,3)\}$ est une base de $\ker (f-2Id)$ et que $\{u_3=(3,4,8)\}$ est une base de $\ker (f-3Id)$.
Donc $\dim \ker (f+Id)=\dim \ker (f-2Id)= \dim \ker (f-3Id)=1$.
\vskip0mm
On vérifie que ${\cal B}' = \{u_1,u_2,u_3\}$ est une base de $\Rr^3$ et que la matrice de $f$ dans cette base
est :  $B= \left(\begin{array}{rcc}-1&0&0\\0&2&0\\ 0&0&3\end{array}\right).$
\vskip0mm
Soit $g= (f+Id)o(f-2Id)o(f-3Id)$. On pose $\lambda_1=-1, \lambda_2=2, \lambda_3=3 $. On vérifie que pour $i=1,2,3$, $g(u_i)=(\lambda_i-\lambda_1)(\lambda_i-\lambda_2)(\lambda_i-\lambda_3)u_i=0$. Par conséquent, $g$ est l'application nulle, puisque ${\cal B}' = \{u_1,u_2,u_3\}$ est une base de $\Rr^3$.
\end{explanations}
\end{question}
}
